\documentclass[../main.tex]{subfiles}
\begin{document}

\chapter{Discusión}
\label{C:discusion}

\section{Síntesis de los hallazgos principales}
\label{sec:sintesis-hallazgos}

El experimento produce tres hallazgos principales. El primero es el rechazo de H1:
las diferencias de \gls{mttr} entre las estrategias~B y~C no son estadísticamente
significativas en los escenarios donde ambas pueden recuperarse (E1, E3, E4). El
segundo es la confirmación de H3: ante la destrucción activa del repositorio (E2),
la estrategia~B falla en el 100\,\% de los casos mientras que C mantiene
disponibilidad total ($p < 0.001$). El tercero es la confirmación de H2: los
tiempos de detección mediante \gls{siem} (mediana 5.7--41.9~s) son cinco órdenes
de magnitud inferiores a los reportados para la detección manual.

Estos tres hallazgos, tomados en conjunto, forman un argumento coherente para la
adopción de una estrategia de defensa en profundidad: el almacenamiento inmutable
garantiza que siempre habrá algo que restaurar, mientras que el \gls{siem}
garantiza que el tiempo entre el inicio del ataque y la respuesta del equipo de
seguridad se mide en segundos, no en meses.

\section{Interpretación del hallazgo principal}

\subsection{Por qué el rechazo de H1 no invalida Object Lock}

El resultado más importante del experimento no es la diferencia de \gls{mttr}
entre estrategias B y C, sino la diferencia de disponibilidad ante un atacante que
destruye activamente las copias. En el escenario~E2, que simula el comportamiento
documentado de grupos como \emph{Ryuk} o \emph{BlackCat}, la estrategia~B falla
en el 100\,\% de los casos mientras que la estrategia~C mantiene una tasa de
éxito del 100\,\%.

El rechazo de H1 en su formulación original no invalida la utilidad práctica de
Object Lock. La hipótesis planteaba una reducción de \gls{mttr} de al menos el
30\,\%, un criterio que asume implícitamente que ambas estrategias son igualmente
capaces de recuperarse. Cuando se incluye el escenario~E2, la comparación más
correcta no es entre tiempos de recuperación sino entre \emph{tasas de
recuperabilidad}.

Esta reformulación tiene implicaciones directas para la toma de decisiones en
organizaciones: la justificación del coste adicional de Object Lock no debe basarse
en la promesa de recuperaciones más rápidas, sino en la garantía de que la
recuperación sea posible en absoluto cuando el atacante ha comprometido las
credenciales de administración.

\subsection{La dicotomía recuperable vs. irrecuperable}

La distinción entre estrategias B y C no es cuantitativa sino cualitativa: no se
trata de recuperar en 10~s vs. 9~s, sino de recuperar vs. no poder recuperar. Esta
dicotomía binaria tiene consecuencias directas sobre el modelo de decisión que
deben usar las organizaciones:

\begin{itemize}
  \item Con la estrategia~B, la pregunta relevante es «¿cuándo fue la última vez
  que el atacante accedió al servidor de backup?». Si el atacante lleva semanas
  en la red (como es habitual según IBM X-Force~2024 \cite{ibm2024}), el
  repositorio puede haber sido destruido silenciosamente antes del cifrado.

  \item Con la estrategia~C, la pregunta es irrelevante: el repositorio existe
  y está íntegro independientemente de lo que el atacante haya hecho con las
  credenciales, siempre que no hayan expirado los objetos por fin de retención.
\end{itemize}

\section{El escenario E2 como evidencia experimental central}

La celda E2$\times$B (0/10 éxitos) frente a E2$\times$C (10/10 éxitos) constituye
la evidencia experimental más sólida del trabajo. La destrucción del repositorio
SFTP es técnicamente sencilla: con acceso \gls{ssh} y las credenciales del usuario
\texttt{restic-b}, basta ejecutar \texttt{rm -rf /srv/backups/restic-b/*} para
dejar a la organización sin posibilidad de recuperación. Esta operación tarda menos
de un segundo y es irreversible.

Object Lock COMPLIANCE bloquea la operación equivalente sobre el \emph{bucket} S3
con un error \texttt{AccessDenied} independientemente de las credenciales
presentadas. No hay mecanismo de \emph{bypass}: ni el usuario raíz de MinIO ni el
administrador de la plataforma pueden borrar un objeto en retención COMPLIANCE
antes de que expire el plazo.

Esto no significa que la estrategia~C sea completamente invulnerable. Un atacante
con acceso físico al servidor donde corre MinIO, o con capacidad de modificar el
binario, podría eludir la protección. Sin embargo, este vector requiere un nivel
de compromiso de la infraestructura que va más allá del modelo de amenaza del
ransomware convencional. En un entorno cloud (Amazon S3), el atacante necesitaría
comprometer la propia infraestructura de Amazon, lo que está fuera del alcance de
la práctica totalidad de los actores de amenaza conocidos.

\section{Sobre la estrategia de detección con FIM y runner.jsonl}

\subsection{Justificación de la alternativa a auditd}

La incompatibilidad de \texttt{auditd} con el \emph{kernel} CachyOS no es un
defecto sino una característica de seguridad: los parches de endurecimiento
restringen el acceso al subsistema de auditoría desde contenedores precisamente
para evitar interferencias maliciosas. La solución adoptada (\gls{fim} más
\texttt{runner.jsonl}) es representativa de los despliegues de Wazuh en producción
real, donde \texttt{auditd} raramente se activa por su \emph{overhead}.

La validación en máquina virtual Ubuntu~24.04 (sección~\ref{sec:vm-reproducibilidad})
confirmó que en entornos con \emph{kernel} estándar, \texttt{auditd} sí está
disponible desde contenedores privilegiados, lo que podría complementar la detección
basada en \gls{fim} con instrumentación a nivel de \emph{syscall}.

\subsection{Ventajas del enfoque basado en eventos estructurados}
\label{sec:ventajas-runner}

Los eventos \gls{json} del \emph{runner} tienen semántica de alto nivel
directamente interpretable
(\mbox{\texttt{immutability\_verified}},
\mbox{\texttt{lateral\_jump}}),
generan un volumen manejable (5--15 eventos por ejecución frente a los miles que
generaría \texttt{auditd}) y son completamente portables entre \emph{kernels}.

Adicionalmente, los eventos son semánticamente significativos desde el punto de
vista de la seguridad: el analista de seguridad recibe eventos de negocio que
describen directamente lo que está ocurriendo, sin necesidad de interpretar
secuencias de llamadas al sistema de bajo nivel. Esto reduce el tiempo de respuesta
y facilita la automatización de la respuesta a incidentes.

\subsection{Limitaciones de la detección basada en FIM}

La principal limitación del \gls{fim} para la detección de ransomware es la
dependencia de umbrales de frecuencia. Las reglas 100200 y 100202 requieren que
se detecten al menos 15 y 10 cambios en directorios monitorizados antes de generar
una alerta. Un cifrador extremadamente rápido podría completar el cifrado antes de
superar esos umbrales.

En el experimento, las reglas del \emph{runner} (100205, 100206) disparan antes
que el \gls{fim} en la mayoría de las ejecuciones de E1, precisamente porque el
evento \texttt{execute} llega al \emph{indexer} antes de que el \gls{fim} haya
acumulado 15~cambios. Esto indica que en un despliegue real, complementar el
\gls{fim} con eventos de aplicación es una estrategia más robusta que depender
exclusivamente del \gls{fim}.

La calibración de los umbrales de frecuencia requiere análisis del perfil de uso
normal del directorio monitorizado. Un umbral demasiado bajo genera falsos positivos
(cualquier operación masiva de ficheros legítima dispara la alerta); un umbral
demasiado alto puede no detectar cifradores muy rápidos.

\section{Sensibilidad de los resultados a la calibración de umbrales}
\label{sec:sensibilidad}

Un aspecto que no se ha discutido hasta ahora es en qué medida los resultados
presentados dependen de los valores concretos de \texttt{frequency} y
\texttt{timeframe} elegidos para las reglas \gls{fim} descritas en la
sección~\ref{sec:reglas} del capítulo~\ref{C:disenyo}, y no solo de la
arquitectura general del sistema de detección.

Los umbrales actuales (por ejemplo, 15~eventos en 120~s para la regla~100200,
o 5~eventos en 30~s para la regla~100211) se calibraron durante la fase~2 del
proyecto descrita en la sección de metodología del capítulo~\ref{C:introduccion},
observando el volumen de eventos \gls{fim} generado tanto por la actividad
legítima del sistema de ficheros como por el cifrador sintético del
escenario~E1. Un umbral significativamente más alto retrasaría la detección
--- con el riesgo de que un cifrador más rápido que el sintético utilizado
complete su trabajo antes de que se supere el umbral, como ya se advierte en la
sección~\ref{sec:ventajas-runner} --- mientras que un umbral más bajo
incrementaría la sensibilidad a costa de un mayor riesgo de falsos positivos
ante operaciones legítimas de copia o sincronización masiva de ficheros.

Es importante señalar que esta sensibilidad afecta únicamente a las reglas
basadas en \gls{fim} (100200, 100202, 100211), no a las basadas en eventos
estructurados del \emph{runner} (100205, 100206, 100212--100214, 100220,
100221, 100230, 100231), que disparan de forma binaria ante la presencia de un
evento concreto y no dependen de ningún parámetro de calibración. Dado que,
como se observa en la sección~\ref{sec:ventajas-runner}, las reglas basadas en
eventos del \emph{runner} son habitualmente las primeras en disparar durante
las ejecuciones de E1 y E2, el resultado cualitativo central del trabajo ---
que el escenario~E2 revela una diferencia de disponibilidad, no de velocidad,
entre estrategias~B y~C --- es en gran medida robusto frente a la calibración
específica de los umbrales de \gls{fim}: incluso si estos se calibraran de
forma distinta, la imposibilidad de recuperación de la estrategia~B ante la
destrucción del repositorio no depende de qué mecanismo la detecte primero,
sino del resultado binario \texttt{recovery\_status} registrado por el
\emph{playbook} de recuperación, completamente independiente de la
instrumentación \gls{siem}. Una caracterización sistemática de esta
sensibilidad --- por ejemplo, mediante un barrido de valores de
\texttt{frequency} y \texttt{timeframe} y su efecto sobre $\tdetect$ --- no se
realizó en este trabajo y se plantea como posible extensión metodológica del
diseño experimental actual.

\section{Comparación con trabajos relacionados}

Los tiempos de detección observados (5.7--41.9~s) son comparables a los reportados
por \textcite{almashhadani2019} para análisis de tráfico de red ($< 30$~s), aunque
las dos aproximaciones son complementarias: el análisis de red detecta el tráfico
de exfiltración durante la transferencia, mientras que el \gls{fim} detecta la
preparación del ataque (creación del \emph{script} en \texttt{/tmp}).

Las tasas de detección son consistentes con las de \textcite{moussaileb2021} para
análisis de acceso a ficheros ($> 95\,\%$). La diferencia fundamental respecto a
los trabajos existentes es la inclusión del escenario de destrucción activa de
\emph{backups} (E2) y la comparación experimental de estrategias de recuperación.

\textcite{cabaj2016} proponen un sistema de mitigación basado en \gls{sdn} que
podría complementar el laboratorio: la detección basada en \gls{fim} podría
alimentar directamente el sistema \gls{sdn} para aislar automáticamente el
segmento de red comprometido en cuanto se detecte el cifrado activo.

La comparación con los tres trabajos metodológicamente próximos introducidos en
el capítulo~\ref{C:estado_arte} --- \textcite{gehani2023}, \textcite{vonderassen2024}
y \textcite{netsecbed2026} --- confirma, tras observar los resultados
experimentales, la complementariedad ya anticipada en la revisión bibliográfica.
El \emph{dataset} reproducible de \textcite{gehani2023} sobre comportamiento
real de binarios de ransomware sería un candidato natural para sustituir el
cifrador sintético del escenario~E1 en la línea de trabajo futuro TF2
(capítulo~\ref{C:conclusiones}), aportando patrones de acceso a ficheros más
representativos sin comprometer la reproducibilidad del laboratorio. La
mitigación activa a nivel de sistema de ficheros de GuardFS~\cite{vonderassen2024}
actuaría, de integrarse, \emph{antes} en la cadena de defensa que el mecanismo
evaluado en este trabajo: mientras que GuardFS intenta reducir la cantidad de
datos perdidos durante el propio cifrado, el par almacenamiento
inmutable~+~\gls{siem} evaluado aquí garantiza la recuperación \emph{después}
de que el cifrado se haya completado, con independencia de cuánto daño se haya
producido. Ambas capas son, por tanto, complementarias y no mutuamente
excluyentes, en línea con el argumento de defensa en profundidad de la
sección~\ref{sec:sintesis-hallazgos} de este mismo capítulo.

\section{Limitaciones del estudio}

\begin{description}
  \item[Entorno sintético] El laboratorio simula una infraestructura empresarial
  mínima con contenedores en una misma máquina. Las latencias de red reales, la
  heterogeneidad del \emph{software} instalado y la presencia de antivirus pueden
  afectar tanto a los tiempos de ataque como a los de detección y recuperación.
  En particular, un servidor de ficheros real con terabytes de datos produciría
  tiempos de cifrado muy superiores a los del laboratorio, lo que podría amplificar
  las diferencias de \gls{mttr} entre escenarios.

  \item[Tamaño de muestra] Diez repeticiones por celda es suficiente para
  estimaciones descriptivas pero limita la potencia estadística. Con muestras
  mayores, diferencias de efecto mediano (Cliff~$\delta \approx 0.36$ en E4)
  podrían alcanzar significación estadística. Para detectar efectos medianos con
  potencia del 80\,\% se necesitarían al menos 30 repeticiones por celda,
  lo que supondría 360 ejecuciones y aproximadamente 15 horas de experimento.

  \item[Escenarios simplificados] Los cuatro escenarios implementados no cubren la
  totalidad de las técnicas de MITRE ATT\&CK. En particular, no se simula la fase
  de reconocimiento previo al ataque ni la persistencia del agente malicioso. En
  ataques reales, el atacante puede permanecer en la red durante semanas o meses
  antes de desplegar el cifrador, lo que permite la destrucción silenciosa y gradual
  de los \emph{snapshots} del repositorio de \emph{backup}.

  \item[Cifrador sintético] El cifrador utilizado en E1 es un \emph{script} Python
  con AES-GCM secuencial que no replica la optimización de los cifradores reales
  como \emph{LockBit}~3.0, que cifra varios gigabytes por segundo mediante E/S
  asíncrona y múltiples hilos.

  \item[Estrategia B sin rotación] El repositorio \gls{sftp} de la estrategia~B
  se inicializa con un único \emph{snapshot}. En producción, los repositorios
  acumulan múltiples \emph{snapshots} distribuidos en el tiempo, lo que haría más
  costoso para el atacante destruir todas las copias.

  \item[Overhead de monitorización no medido] Durante el experimento no se midió
  sistemáticamente el impacto en \gls{cpu} y E/S de disco del \gls{fim} en tiempo
  real sobre el \emph{fileserver} durante el cifrado activo. Esta medición
  permitiría determinar si la monitorización introduce algún sesgo en los tiempos
  observados.
\end{description}

\section{Implicaciones prácticas}
\label{sec:implicaciones-practicas}

Los resultados tienen implicaciones directas para la arquitectura de sistemas de
\emph{backup} en organizaciones:

\begin{enumerate}
  \item \textbf{Almacenamiento inmutable obligatorio}: implementar al menos una
  copia con Object Lock COMPLIANCE en una cuenta con credenciales completamente
  separadas. El coste adicional de MinIO \emph{on-premise} frente a un servidor
  SFTP convencional es mínimo comparado con el riesgo eliminado.

  \item \textbf{SIEM con FIM activo}: desplegar \gls{fim} en tiempo real sobre los
  directorios críticos. Los tiempos observados (mediana 6--8~s en E1) son
  suficientemente cortos para activar procedimientos de contención antes de que el
  cifrado se complete en la mayoría de los escenarios reales.

  \item \textbf{Pruebas periódicas de recuperación}: verificar mensualmente que la
  restauración desde el repositorio inmutable funciona correctamente. La gestión de
  \emph{locks} de \restic con Object Lock requiere \texttt{--no-lock} en todas las
  operaciones.

  \item \textbf{Separación de credenciales}: usar claves \gls{ssh} únicas para
  cada servidor de \emph{backup}, distintas de las de administración, almacenadas
  en un gestor de secretos con acceso auditado.

  \item \textbf{Monitorización del servidor de backup}: desplegar \gls{fim} no
  solo en los servidores víctima sino también en el servidor de \emph{backup},
  para detectar intentos de destrucción del repositorio (E2) antes de que se
  completen.
\end{enumerate}

\section{Análisis coste-beneficio de la adopción de Object Lock}
\label{sec:coste-beneficio}

Las implicaciones prácticas enumeradas en la sección anterior recomiendan la
adopción de almacenamiento inmutable, pero toda recomendación técnica dirigida a
una organización real debe acompañarse de una estimación de su coste. Esta
sección ofrece un análisis \emph{orientativo}, no una cotización, con el único
propósito de contextualizar la magnitud del coste adicional frente al riesgo que
mitiga.

Según la documentación pública de precios de Amazon~S3 vigente en 2026, el
almacenamiento en la clase \emph{S3~Standard} tiene un coste aproximado de
0,023~\$/GB al mes para los primeros 50~TB, y la activación de Object Lock no
añade ningún cargo adicional sobre el precio base de almacenamiento: se paga
exactamente la misma tarifa por GB almacenado, con o sin protección WORM
activa. El coste diferencial real de la estrategia~C frente a la estrategia~B no
proviene, por tanto, del precio por gigabyte, sino de dos factores indirectos:
el hecho de que los objetos protegidos no pueden eliminarse antes de expirar su
periodo de retención --- lo que impide aplicar políticas de purga agresivas que
reducirían el volumen almacenado --- y un incremento en el número de llamadas a
la \gls{api} necesarias para consultar y gestionar los metadatos de retención de
cada objeto.

\begin{table}[!h]
  \centering
  \caption[Estimación orientativa de coste mensual de almacenamiento para un
           repositorio de backup de 1\,TB bajo las estrategias B y C.]
          {Estimación orientativa de coste mensual de almacenamiento para un
           repositorio de backup de 1\,TB bajo las estrategias B y C.}
  \label{tab:coste-beneficio}
  \rowcolors{2}{tabla-par}{white}
  \begin{tabularx}{\textwidth}{X C{3cm} C{3cm}}
    \toprule
    \rowcolor{tabla-cabecera}
    \textcolor{white}{\textbf{Concepto}} &
    \textcolor{white}{\textbf{Estrategia B (SFTP)}} &
    \textcolor{white}{\textbf{Estrategia C (Object Lock)}} \\
    \midrule
    Almacenamiento (1\,TB, tarifa S3 Standard) & $\approx$23,55\,\$/mes & $\approx$23,55\,\$/mes \\
    Sobrecoste directo de Object Lock          & --- & 0,00\,\$ (sin cargo adicional) \\
    Flexibilidad de purga anticipada            & Sí & No durante el periodo de retención \\
    Coste de un incidente E2 (repositorio destruido, sin recuperación) & Pérdida total de datos & N/A --- éxito garantizado \\
    \bottomrule
  \end{tabularx}
\end{table}

La conclusión del análisis es que, en términos puramente de precio por gigabyte
almacenado, la estrategia~C no es más cara que la estrategia~B: la diferencia de
coste relevante no es la de la infraestructura de almacenamiento en sí, sino la
del \emph{riesgo retenido}. El coste medio de un incidente de ransomware sin
posibilidad de recuperación, estimado en más de cuatro millones de dólares según
el informe Veeam~2024~\cite{veeam2024} citado en el
capítulo~\ref{C:estado_arte}, hace que incluso una sobrecarga operativa moderada
en la gestión de \gls{api} y en la rigidez de las políticas de retención resulte
despreciable frente al coste esperado de un escenario equivalente a~E2 sin
protección. Esta conclusión coincide, además, con el endurecimiento de los
requisitos de las aseguradoras de ciberriesgo descrito en la
sección~\ref{sec:ciberseguros} del capítulo~\ref{C:estado_arte}: el mercado
asegurador ha llegado independientemente a la misma valoración de que el coste
marginal del almacenamiento inmutable es bajo en relación con el riesgo que
elimina.

Conviene matizar que este análisis no incluye los costes de migración desde una
estrategia SFTP existente, la formación del personal de operaciones en la
gestión de \emph{locks} de \restic con Object Lock (sección~\ref{sec:limitaciones}),
ni el coste de una posible implementación \emph{on-premise} con MinIO --- la
utilizada en este laboratorio --- frente al uso de Amazon~S3 gestionado, que
sustituye el coste de licencia por el coste de hardware propio y su
mantenimiento. Una estimación completa de estos costes queda fuera del alcance
de este trabajo y se plantea como línea de trabajo futuro en el
capítulo~\ref{C:conclusiones}.

\section{Notificación en tiempo real mediante Telegram}
\label{sec:telegram}

Durante el desarrollo del laboratorio se implementó un sistema de notificación en
tiempo real mediante un bot de Telegram integrado con Wazuh. El \emph{script}
\texttt{custom-telegram.py}, desplegado en
\texttt{/var/ossec/\allowbreak{}integrations/} del \emph{manager}, recibe las
alertas de nivel~$\geq 10$ y las reenvía a un grupo privado de Telegram en formato
estructurado. La integración se configura en el fichero \texttt{ossec.conf} del
manager como una integración de tipo \texttt{custom}.

Los siguientes mensajes se generaron durante el experimento:

\begin{lstlisting}[style=bash, caption={Ejemplos de notificaciones Telegram
  generadas durante las ejecuciones del experimento.}, label={lst:telegram}]
Alerta Wazuh -- Nivel 10
Regla: 100200 -- E1: cambio masivo de ficheros en /opt/labprueba
Agente: fileserver | Grupos: lab, e1_local
Timestamp: 2026-07-27T11:20:35.282+0000

Alerta Wazuh -- Nivel 12
Regla: 100240 -- E4: exfiltracion completada
Agente: wazuh-manager | Grupos: lab, e4_exfil
Timestamp: 2026-07-24T19:18:41.536+0000
\end{lstlisting}

Esta integración resultó especialmente útil durante las cuatro horas de ejecución
del experimento, permitiendo detectar inmediatamente cualquier anomalía sin
supervisión continua del terminal. En un entorno de producción, la integración con
Telegram (o con sistemas de mensajería empresarial como Slack o Microsoft Teams)
reduciría significativamente el tiempo de respuesta del equipo de seguridad al
no requerir que el analista esté activamente monitorizando el dashboard.

\section{Validación de reproducibilidad en máquina virtual}
\label{sec:vm-reproducibilidad}

Con el objetivo de validar la reproducibilidad del laboratorio en un entorno
diferente al de desarrollo, se realizó una instalación completa en una máquina
virtual Ubuntu~24.04~LTS desde cero. El proceso siguió exactamente las instrucciones
del Anexo~C sin ninguna modificación previa, instalando únicamente las herramientas
mínimas necesarias:

\begin{lstlisting}[style=bash, caption={Instalación mínima en Ubuntu 24.04 para
  reproducir el laboratorio desde cero.}, label={lst:vm-setup}]
# Actualizar el sistema
sudo apt update && sudo apt upgrade -y

# Instalar Docker Engine y Compose v2
sudo apt install -y docker.io docker-compose-plugin

# Instalar dependencias de Python para el análisis estadístico
sudo apt install -y python3-pip
pip3 install scipy numpy pandas matplotlib --break-system-packages

# Requisito del kernel para OpenSearch
sudo sysctl -w vm.max_map_count=262144
echo "vm.max_map_count=262144" | sudo tee -a /etc/sysctl.conf

# Clonar el repositorio y arrancar
git clone https://github.com/AlejandroMM2/Ransomware-simulator
cd Ransomware-simulator
cp .env.example .env
./generate_certs.sh
docker compose up -d wazuh-indexer
# (esperar ~60s y ejecutar securityadmin)
docker compose up -d
bash scripts/post_start.sh
\end{lstlisting}

El \emph{stack} arrancó correctamente en la máquina virtual con \emph{kernel}
estándar de Ubuntu. La diferencia más notable respecto al entorno de desarrollo
es que en Ubuntu, \texttt{auditd} sí estuvo disponible desde los contenedores
privilegiados, lo que podría complementar la detección basada en \gls{fim} con
instrumentación a nivel de \emph{syscall}.

La ejecución de las reglas de detección y el flujo completo del experimento
funcionaron de forma idéntica al entorno original, confirmando la portabilidad y
reproducibilidad del laboratorio documentada en el Anexo~C. Esta validación es
especialmente relevante porque demuestra que el laboratorio no está vinculado al
\emph{hardware} ni al sistema operativo del equipo de desarrollo, y puede ejecutarse
en cualquier máquina Linux moderna con Docker instalado.

\end{document}
